\section{Methods}\label{section:methods}

\subsection{Methods used in Part A}

One-Qubit bases

\begin{equation}\label{single-qubits}
    \ket{0} = \begin{pmatrix}
        1\\0
    \end{pmatrix}\quad
    \ket{1} = \begin{pmatrix}
        0\\1
    \end{pmatrix}
\end{equation}

Pauli Matrices

\begin{equation}\label{pauli-matrices}
    \sigma_x=\begin{pmatrix}
        0&1\\1&0
    \end{pmatrix}\quad
    \sigma_y=\begin{pmatrix}
        0&-i\\i&0
    \end{pmatrix}\quad
    \sigma_z=\begin{pmatrix}
        1&0\\0&-1
    \end{pmatrix}
\end{equation}

Hadamard and Phase Gates

\begin{equation}
    H = \frac{1}{\sqrt{2}}\begin{pmatrix}
        1&1\\1&-1
    \end{pmatrix}\quad
    S = \begin{pmatrix}
        1&0\\0&-i
    \end{pmatrix}
\end{equation}

Kronecker products for basis states

\begin{equation}
\ket{00} = \ket{0} \otimes \ket{0}
\quad
\ket{01} = \ket{0} \otimes \ket{1}
\end{equation}

\begin{equation}
\ket{10} = \ket{1} \otimes \ket{0}
\quad
\ket{11} = \ket{1} \otimes \ket{1}
\end{equation}

Bell states

\begin{equation}
\ket{\Phi^+} = \frac{\ket{00} + \ket{11}}{\sqrt{2}}
\quad
\ket{\Phi^-} = \frac{\ket{00} - \ket{11}}{\sqrt{2}}
\end{equation}

\begin{equation}
\ket{\Psi^+} = \frac{\ket{01} + \ket{10}}{\sqrt{2}}
\quad
\ket{\Psi^-} = \frac{\ket{01} - \ket{10}}{\sqrt{2}}
\end{equation}

\begin{equation}
    \text{CNOT} = \begin{pmatrix}
        1&0&0&0\\
        0&1&0&0\\
        0&0&0&1\\
        0&0&1&0
    \end{pmatrix}
\end{equation}



\subsection{Methods used in Part B}

We define a symmetric matrix $H\in\mathbb{R}^{2\times2}$

\begin{equation}
    H=\begin{bmatrix}
        H_{11} & H_{12} \\
        H_{21} & H_{22}
    \end{bmatrix}
\end{equation}

We let $H=H_0+H_1$ where

\begin{equation}
    H_0=\begin{bmatrix}
        E_1 & 0 \\
        0 & E_2
    \end{bmatrix}
\end{equation}

is a diagonal matrix.

Similarly,

\begin{equation}
    H_1=\begin{bmatrix}
        V_{11} & V_{12} \\
        V_{21} & V_{22}
    \end{bmatrix}
\end{equation}

Where $V_{ij}$ represents various interaction matrix elements. We can view $H_0$ as the non-interacting solution

\begin{equation}
    H_0\ket{0} = E_1\ket{0}
\end{equation}

\begin{equation}
    H_0\ket{1}=E_2\ket{1}
\end{equation}

We rewrite $H$ with $H_0$ and $H_1$ via Pauli matrices

\begin{equation}
    H_0=\mathcal{E}I+\Omega\sigma_Z
\end{equation}

\begin{equation}
    \mathcal{E} = \frac{E_1+E_2}{2}
\end{equation}

\begin{equation}
    \Omega=\frac{E_1-E_2}{2}
\end{equation}

\begin{equation}
    H_I=c\textbf{I}+\omega_z\sigma_z+\omega_x\sigma_x
\end{equation}

\begin{equation}
    c=\frac{V_{11}+V_{22}}{2}
\end{equation}

\begin{equation}
    \omega_z=\frac{V_{11}-V_{22}}{2}
\end{equation}

\begin{equation}
    \omega_x=V_{12}=V_{21}
\end{equation}

We let our Hamiltonian depend linearly on a strenght parameter $\lambda$

\begin{equation}
    H=H_0+\lambda H_1
\end{equation}

with $\lambda\in[0,1]$ and limits $\lambda=0$, $\lambda=1$.

\subsection{Methods used in Part C}

\subsection{Methods used in Part D}

\subsection{Methods used in Part E}

\subsection{Methods used in Part F}

$H=H_0+H_1+H_2$

\begin{equation}
    H_0=\frac{1}{2}\varepsilon\sum_{p\sigma}\sigma a_{p\sigma}^{\dagger}a_{p\sigma}
\end{equation}

\begin{equation}
    H_1=\frac{1}{2}V\sum_{p,p',\sigma}\sigma a_{p\sigma}^{\dagger}a_{p\sigma}
\end{equation}

\begin{equation}
    H_2=\frac{1}{2}W\sum_{p,p',\sigma}a_{p\sigma}^{\dagger}a_{p'\sigma}^{\dagger}a_{p'-\sigma}a_{p-\sigma}
\end{equation}

\begin{equation}
    H_0=\varepsilon J_z
\end{equation}

\begin{equation}
    H_1=\frac{1}{2}V(J_+^2+J_-^2)
\end{equation}

\begin{equation}
    H_2=\frac{1}{2}W(-N+J_+J_-+J_-J_+)
\end{equation}

\begin{equation}
    H_{J=1} = \begin{pmatrix}
        -\epsilon & 0 & -V \\
        0 & 0 & 0 \\
        -V & 0 & \epsilon
    \end{pmatrix}
\end{equation}

For $J=2$ we have a $5\times5$ matrix given by

\begin{equation}
    H_{J=2}=\begin{pmatrix}
        -2\epsilon & 0 & \sqrt{6}V & 0 & 0 \\
        0 & -\epsilon + 3W  & 0 & 3V & 0 \\
        \sqrt{6}V & 0 & 4W & 0 & \sqrt{6}V \\
        0 & 3V & 0 & \epsilon+3W & 0 \\
        0 & 0& \sqrt{6}V & 0 & 2\varepsilon
    \end{pmatrix}
\end{equation}

\subsection{Methods used in Part G}